\documentclass[11pt]{article}
\usepackage{mystyle}

\title{Rosen --- \S 9.5: Equivalence Relations\\ \large Selected Exercises}
\author{Alston}
\date{Semester 2, 2026}

\begin{document}
\maketitle

% ======================================================================
\begin{exercise}{9}
    Suppose that $A$ is a nonempty set, and $f$ is a function that has
    $A$ as its domain. Let $R$ be the relation on $A$ consisting of
    all ordered pairs $(x,y)$ such that $f(x) = f(y)$.
    \begin{enumerate}[label=(\alph*)]
        \item Show that $R$ is an equivalence relation on $A$.
        \item What are the equivalence classes of $R$?
    \end{enumerate}
\end{exercise}

% ======================================================================
\begin{exercise}{10}
    Suppose that $A$ is a nonempty set and $R$ is an equivalence
    relation on $A$. Show that there is a function $f$ with $A$ as its
    domain such that $(x,y) \in R$ if and only if $f(x) = f(y)$.
\end{exercise}

% ======================================================================
\begin{exercise}{15}
    Let $R$ be the relation on the set of ordered pairs of positive
    integers such that $((a,b),(c,d)) \in R$ if and only if
    $a + d = b + c$. Show that $R$ is an equivalence relation.
\end{exercise}

% ======================================================================
\begin{exercise}{16}
    Let $R$ be the relation on the set of ordered pairs of positive
    integers such that $((a,b),(c,d)) \in R$ if and only if
    $ad = bc$. Show that $R$ is an equivalence relation.
\end{exercise}

% ======================================================================
\begin{exercise}{25}
    Show that the relation $R$ on the set of all bit strings such that
    $s\,R\,t$ if and only if $s$ and $t$ contain the same number of
    1s is an equivalence relation.
\end{exercise}

% ======================================================================
\begin{exercise}{54}
    \textit{A partition $P_1$ is called a \emph{refinement} of the
    partition $P_2$ if every set in $P_1$ is a subset of one of the
    sets in $P_2$.}

    Suppose that $R_1$ and $R_2$ are equivalence relations on a set
    $A$. Let $P_1$ and $P_2$ be the partitions that correspond to
    $R_1$ and $R_2$, respectively. Show that $R_1 \subseteq R_2$ if
    and only if $P_1$ is a refinement of $P_2$.
\end{exercise}

% ======================================================================
\begin{exercise}{56}
    Suppose that $R_1$ and $R_2$ are equivalence relations on the set
    $S$. Determine whether each of these combinations of $R_1$ and
    $R_2$ must be an equivalence relation.
    \begin{enumerate}[label=(\alph*)]
        \item $R_1 \cup R_2$
        \item $R_1 \cap R_2$
        \item $R_1 \oplus R_2$
    \end{enumerate}
\end{exercise}

% ======================================================================
\begin{exercise}{58}
    Each bead on a bracelet with three beads is either red, white, or
    blue. Define the relation $R$ between bracelets as: $(B_1, B_2)$,
    where $B_1$ and $B_2$ are bracelets, belongs to $R$ if and only if
    $B_2$ can be obtained from $B_1$ by rotating it or rotating it and
    then reflecting it.
    \begin{enumerate}[label=(\alph*)]
        \item Show that $R$ is an equivalence relation.
        \item What are the equivalence classes of $R$?
    \end{enumerate}
\end{exercise}

% ======================================================================
\begin{exercise}{60}
    \begin{enumerate}[label=(\alph*)]
        \item Let $R$ be the relation on the set of functions from
        $\mathbb{Z}^+$ to $\mathbb{Z}^+$ such that $(f,g)$ belongs to
        $R$ if and only if $f$ is $\Theta(g)$ (see Section 3.2). Show
        that $R$ is an equivalence relation.
        \item Describe the equivalence class containing $f(n) = n^2$
        for the equivalence relation of part (a).
    \end{enumerate}
\end{exercise}

% ======================================================================
\begin{exercise}{63}
    Do we necessarily get an equivalence relation when we form the
    transitive closure of the symmetric closure of the reflexive
    closure of a relation?
\end{exercise}

% ======================================================================
\begin{exercise}{68}
    Let $p(n)$ denote the number of different equivalence relations on
    a set with $n$ elements (and, by Theorem 2, the number of
    partitions of a set with $n$ elements). Show that $p(n)$ satisfies
    the recurrence relation
    \[
    p(n) = \sum_{j=0}^{n-1} \binom{n-1}{j} p(n-j-1)
    \]
    and the initial condition $p(0) = 1$. (\emph{Note}: The numbers
    $p(n)$ are called \emph{Bell numbers} after the American
    mathematician E. T. Bell.)
\end{exercise}

\end{document}